\chapter{Logarithme, puissances et modèles multiplicatifs}
\section*{Exercice 16.1 -- Domaine et produit}
[Base] Résoudre $\ln(x-2)+\ln x=\ln3$.

\section*{Exercice 16.2 -- Seuil}
[Méthode] Trouver le plus petit entier $n$ avec $0{,}8^n<0{,}1$.

\section*{Exercice 16.3 -- Limite en zéro}
[Base] Calculer $\lim_{x\to0^+}x^2\ln x$.

\section*{Exercice 16.4 -- Puissance réelle}
[Approfondissement] Dériver $x^{\sqrt2}$ sur $]0,+\infty[$ en revenant à sa définition.

\section*{Exercice 16.5 -- Logarithme d’une somme}
[Critique] Résoudre $\ln(x+1)=\ln x+\ln2$ sans utiliser de règle fausse.

\section*{Exercice 16.6 -- Une limite exponentielle}
[Preuve] Montrer $(1+2/n)^n\to e^2$.
